\documentclass{article}

\usepackage{amsmath,amssymb,amsthm}
\usepackage{booktabs}
\usepackage{graphicx}
\usepackage{hyperref}
\usepackage{microtype}
\usepackage{xcolor}
\usepackage{cleveref}

\title{Contour Certificates for Stable Learned Dynamics:\\
A Rouch\'e--Perturbation Framework for State-Space Models}
\author{Author Names Omitted}
\date{}

\newtheorem{theorem}{Theorem}
\newtheorem{proposition}{Proposition}
\newtheorem{lemma}{Lemma}
\newtheorem{corollary}{Corollary}
\newtheorem{remark}{Remark}
\theoremstyle{definition}
\newtheorem{definition}{Definition}

\newcommand{\C}{\mathbb{C}}
\newcommand{\R}{\mathbb{R}}
\newcommand{\N}{\mathbb{N}}
\newcommand{\rho}{\varrho}       % spectral radius
\newcommand{\norm}[1]{\left\|#1\right\|}
\newcommand{\abs}[1]{\left|#1\right|}

% placeholder for numerical results not yet filled in
\newcommand{\TBD}[1]{\textcolor{red}{[TBD: #1]}}

\begin{document}
\maketitle

\begin{abstract}
Learned state-space models rely on stable transition dynamics for reliable
long-horizon prediction and gradient-stable training.  We develop
contour-based \emph{sufficient} certificates for Schur stability of learned
transition matrices using scalar and matrix-valued Rouch\'e-type perturbation
arguments.  The resulting determinant and resolvent certificates provide
computable post-hoc stability tests anchored to known-stable reference
systems.  We derive a deterministic finite-grid verification theorem that
converts sampled margin checks into continuum guarantees, and study
reference-selection algorithms that improve certificate margins over naive
scalar baselines.  Synthetic experiments characterise the conservatism,
numerical conditioning, and wall-clock scalability of the certificates across
normal, non-normal, near-defective, block-diagonal, and structured
state-space transition matrices.  Our central finding is that resolvent
certificates outperform determinant certificates across all matrix families,
and that reference selection is the single largest lever for reducing false
negatives among stable matrices.
\end{abstract}

%% ============================================================
\section{Introduction}
\label{sec:intro}

Discrete-time state-space sequence models (SSMs) maintain a hidden state
$h_t \in \C^n$ updated by
\begin{equation}
  h_t = A h_{t-1} + B u_t, \qquad y_t = C h_t + D u_t,
  \label{eq:ssm}
\end{equation}
where $A \in \C^{n \times n}$ is the transition matrix, $u_t$ is the input,
and $y_t$ is the output.  The stability of long-horizon behaviour and the
well-conditioning of gradients through time depend on the transition matrix
satisfying Schur stability: every eigenvalue of $A$ must lie strictly inside
the open unit disk, i.e.\ $\varrho(A) < 1$ where $\varrho$ denotes the
spectral radius.

Contemporary SSM architectures~\cite{gu2022s4,smith2023lru,gu2023mamba} either
constrain $A$ to a stable parameterisation by design (diagonal, HIPPO, or
unitary) or impose spectral penalties during training.  Direct spectral-radius
tests after training are exact but offer limited interpretability: knowing that
$\varrho(A) = 0.97$ does not reveal which components of $A$ are closest to
instability, nor how much perturbation the system can absorb while remaining
stable.  Lyapunov- and $\mu$-analysis methods provide richer certificates but
require solving semidefinite or structured-singular-value problems that scale
poorly with $n$.

This paper asks a different question: \emph{when do Rouch\'e-type
perturbation arguments around known-stable reference dynamics yield useful
sufficient certificates for Schur stability?}  Classical Rouch\'e's theorem
is a one-line complex-analytic result that bounds the number of zeros of a
holomorphic function inside a contour by comparing it to a reference function
whose zeros are known.  Applied to characteristic polynomials or to
matrix-valued resolvents on the unit circle, it converts the question ``is
$A$ Schur stable?'' into ``is the perturbation $A - A_0$ small, in an
appropriate norm, relative to the reference dynamics $A_0$?''

Our contributions are:
\begin{enumerate}
  \item A \emph{scalar determinant certificate} (Theorem~\ref{thm:det})
    applying classical Rouch\'e directly to characteristic polynomials.
  \item A \emph{matrix-valued resolvent certificate} (Theorem~\ref{thm:res})
    that avoids ill-conditioning of high-degree polynomials and connects
    naturally to the stability radius and small-gain framework.
  \item A \emph{deterministic finite-grid verification theorem}
    (Theorem~\ref{thm:grid}) that upgrades a sampled margin check to a
    continuum certificate using a Lipschitz bound on the margin function.
  \item \emph{Reference-selection algorithms} (Section~\ref{sec:refsel}) that
    maximise certificate margins over structured reference families, reducing
    false-negative rates substantially over scalar baselines.
  \item \emph{Synthetic experiments}
    (Section~\ref{sec:experiments}) characterising conservatism, false-negative
    rates, numerical conditioning, and wall-clock scalability.
\end{enumerate}

We emphasise throughout that the certificates are \emph{sufficient but not
necessary}: a matrix can be Schur stable and yet fail both certificates if no
close stable reference exists.  The goal is not to replace spectral-radius
tests but to provide a complementary tool that is interpretable, structured,
and amenable to post-hoc verification.

%% ============================================================
\section{Background}
\label{sec:background}

\subsection{Schur stability}

\begin{definition}
A matrix $A \in \C^{n \times n}$ is \emph{Schur stable} if
$\varrho(A) := \max\{\abs{\lambda} : \lambda \in \sigma(A)\} < 1$.
\end{definition}

Schur stability is the discrete-time analogue of Hurwitz stability and is
equivalent to requiring that the characteristic polynomial
$p_A(z) = \det(zI - A)$ has all roots strictly inside the open unit disk.
For a single matrix this can be checked exactly in $O(n^3)$ time by computing
eigenvalues, or by the Schur--Cohn criterion~\cite{jury1964inners}.

\subsection{Rouch\'e's theorem}

\begin{theorem}[Rouch\'e~\cite{ahlfors1979complex}]
\label{thm:rouche_classical}
Let $f$ and $g$ be holomorphic inside and on a simple closed contour $C$.
If $\abs{f(z) - g(z)} < \abs{f(z)}$ for all $z \in C$, then $f$ and $g$
have the same number of zeros inside $C$, counted with multiplicity.
\end{theorem}

The standard proof uses the argument principle: the winding number of $g/f$
around the origin equals the difference in zero counts, and the strict
domination hypothesis forces that winding number to zero.  An equivalent
statement replaces the strict inequality with $\abs{g(z)} > 0$, i.e.\
$f(z) \neq 0$ on $C$, which is automatic from the hypothesis.

\subsection{Pseudospectra and the stability radius}

The \emph{$\varepsilon$-pseudospectrum} of $A$ is
$\Lambda_\varepsilon(A) = \{z \in \C : \norm{(zI-A)^{-1}} \geq
\varepsilon^{-1}\}$.  The \emph{complex stability radius} of $A$ with
respect to perturbations in $\C^{n\times n}$ is
\[
  r_{\C}(A) = \min\{\norm{\Delta} : \varrho(A + \Delta) \geq 1\}
  = \min_{|z|=1} \sigma_{\min}(zI - A),
\]
where $\sigma_{\min}$ denotes the smallest singular value~\cite{hinrichsen1986stability,trefethen2005spectra}.
The resolvent certificate of Section~\ref{sec:res} is directly related to
this quantity: if $\norm{(zI-A_0)^{-1}(A_0 - A)} < 1$ on the unit circle,
then $A$ is a perturbation of $A_0$ that lies inside the stability ball of
$A_0$ under the resolvent norm.

\subsection{SSMs and stability constraints}

Modern SSM layers such as S4~\cite{gu2022s4}, LRU~\cite{smith2023lru}, and
Mamba~\cite{gu2023mamba} typically restrict $A$ to a diagonal or diagonal-plus-low-rank
structure with eigenvalues parameterised in polar form to facilitate
training-time stability constraints.  Our certificates apply to any $A$
post-training and are most informative when $A$ is close to a known-stable
reference.

%% ============================================================
\section{Scalar Determinant Certificate}
\label{sec:det}

Let $p_A(z) = \det(zI - A)$ and $p_0(z) = \det(zI - A_0)$ be the
characteristic polynomials of $A$ and a reference $A_0$.  Define the
\emph{determinant margin} at a point $z \in \C$ as
\[
  m_{\det}(z) = \abs{p_0(z)} - \abs{p_A(z) - p_0(z)}.
\]

\begin{theorem}[Scalar contour certificate]
\label{thm:det}
Suppose $A_0$ is Schur stable.  If $m_{\det}(z) > 0$ for every $z$ on the
unit circle $|z| = 1$, then $A$ is Schur stable.
\end{theorem}

\begin{proof}
Apply Theorem~\ref{thm:rouche_classical} with $f = p_0$, $g = p_A$, and
$C = \{z : |z| = 1\}$.  The hypothesis $\abs{p_0(z)} > \abs{p_A(z) - p_0(z)}$
for all $|z| = 1$ is precisely the Rouch\'e condition.  Since $A_0$ is Schur
stable, $p_0$ has exactly $n$ zeros strictly inside the unit disk (all $n$
eigenvalues of $A_0$) and none on the unit circle.  Rouch\'e's theorem
therefore implies that $p_A$ also has $n$ zeros inside the unit disk.
Because $p_A$ is a monic polynomial of degree $n$, all $n$ roots are
accounted for, so every eigenvalue of $A$ lies inside the open unit disk,
i.e.\ $A$ is Schur stable.
\end{proof}

\begin{remark}
For $n = 1$ the certificate reduces to $\abs{1 - a_0} > \abs{a - a_0}$ on
$|z| = 1$, which is $|a - a_0| < \min_{|z|=1}|z - a_0| = 1 - |a_0|$, i.e.\
the open ball around $a_0$ of radius $1 - |a_0|$.  This is tight.
\end{remark}

\begin{remark}
For large $n$, the determinant $\abs{p_0(z)}$ can be exponentially large or
small, introducing severe numerical cancellation.  This motivates the
matrix-valued certificate in the next section.
\end{remark}

%% ============================================================
\section{Matrix-Valued Resolvent Certificate}
\label{sec:res}

Define the \emph{resolvent quantity} and \emph{resolvent margin} at $z$ as
\begin{align}
  q_{\mathrm{res}}(z) &= \norm{(zI - A_0)^{-1}(A_0 - A)}, \\
  m_{\mathrm{res}}(z) &= 1 - q_{\mathrm{res}}(z).
\end{align}
Here $\norm{\cdot}$ denotes the matrix 2-norm (largest singular value) unless
otherwise stated.

\begin{theorem}[Resolvent contour certificate]
\label{thm:res}
Suppose $A_0$ is Schur stable and
\[
  \sup_{|z|=1} \norm{(zI - A_0)^{-1}(A_0 - A)} < 1.
\]
Then $A$ is Schur stable.
\end{theorem}

\begin{proof}
Let $F_0(z) = zI - A_0$ and $F(z) = zI - A$.  Note that $F(z) - F_0(z) =
A_0 - A$.  For any $|z| = 1$, write
\[
  F(z) = F_0(z)\bigl(I + F_0(z)^{-1}(A_0 - A)\bigr).
\]
The hypothesis implies $\norm{F_0(z)^{-1}(A_0 - A)} < 1$, so by the
Neumann series the factor $I + F_0(z)^{-1}(A_0 - A)$ is invertible, and
therefore $\det F(z) = \det(zI - A) \neq 0$ for every $|z| = 1$.

It remains to show all $n$ eigenvalues lie \emph{inside} (not merely off) the
unit circle.  Consider the homotopy $A_t = A_0 + t(A - A_0)$ for
$t \in [0, 1]$.  For each $t$:
\[
  \norm{(zI - A_0)^{-1}(A_0 - A_t)} = t \norm{(zI - A_0)^{-1}(A_0 - A)} < 1,
\]
so $\det(zI - A_t) \neq 0$ on the unit circle for all $t$.  The number of
eigenvalues of $A_t$ inside the unit disk is therefore constant in $t$ (it is
a continuous integer-valued function of $t$).  At $t = 0$ all $n$ eigenvalues
of $A_0$ are inside the disk; at $t = 1$ all $n$ eigenvalues of $A$ are
inside the disk.  Hence $A$ is Schur stable.
\end{proof}

\begin{remark}[Connection to small-gain]
The condition $\sup_{|z|=1} \norm{(zI-A_0)^{-1}(A_0-A)} < 1$ is the
small-gain condition: the open-loop gain of the feedback system formed by the
reference resolvent $(zI - A_0)^{-1}$ and the perturbation $(A_0 - A)$ is
less than one.  Theorem~\ref{thm:res} can therefore be viewed as a special
case of the small-gain theorem for discrete-time causal operators
evaluated on the unit circle~\cite{desoer1975feedback}.
\end{remark}

\begin{remark}[Diagonal case]
When both $A$ and $A_0$ are diagonal, $A = \mathrm{diag}(a_i)$ and $A_0 =
\mathrm{diag}(a_{0,i})$, the matrix $(zI - A_0)^{-1}(A_0 - A)$ is diagonal
with entries $(a_{0,i} - a_i)/(z - a_{0,i})$.  The certificate reduces to
$\max_i \sup_{|z|=1} |a_{0,i} - a_i|/|z - a_{0,i}| < 1$, i.e.\
$|a_{0,i} - a_i| < 1 - |a_{0,i}|$ for every $i$.
This has a clean geometric interpretation: each diagonal entry of $A$ must
lie in the open disk of radius $1 - |a_{0,i}|$ centred at $a_{0,i}$.
\end{remark}

%% ============================================================
\section{Finite-Grid Verification}
\label{sec:grid}

In practice, the unit circle is approximated by a finite equispaced grid
$\theta_k = 2\pi k / N$, $k = 0, \ldots, N-1$.  We need conditions under
which positivity of the margin on the grid guarantees positivity everywhere.

\begin{theorem}[Deterministic finite-grid certificate]
\label{thm:grid}
Let $m : [0, 2\pi] \to \R$ be a $2\pi$-periodic $L$-Lipschitz function
representing either $m_{\det}$ or $m_{\mathrm{res}}$ as a function of
$\theta = \arg(z)$.  Let $\theta_k = 2\pi k/N$ for an equispaced grid of
$N \geq 4$ points.  If
\[
  \min_{k} m(\theta_k) > \frac{\pi L}{N},
\]
then $m(\theta) > 0$ for all $\theta \in [0, 2\pi]$.
\end{theorem}

\begin{proof}
For any $\theta \in [0, 2\pi]$, let $\theta_k$ be the nearest grid point.
The grid spacing is $2\pi/N$, so $|\theta - \theta_k| \leq \pi/N$.  By the
Lipschitz condition,
\[
  m(\theta) \geq m(\theta_k) - L|\theta - \theta_k| \geq m(\theta_k) - \frac{\pi L}{N}
  > \frac{\pi L}{N} - \frac{\pi L}{N} = 0. \qedhere
\]
\end{proof}

\begin{remark}[Two modes of use]
Theorem~\ref{thm:grid} requires a \emph{validated} upper bound $L$ on the
Lipschitz constant.  In the absence of such a bound, the sampled minimum
$\min_k m(\theta_k) > 0$ is a useful diagnostic (it certifies the matrix
has no eigenvalue ``visible'' from the grid) but does not constitute a
mathematical certificate.  Our experiments report both.
\end{remark}

\paragraph{Lipschitz bound for the resolvent margin.}
The resolvent margin $m_{\mathrm{res}}(\theta) = 1 - \norm{(e^{i\theta}I -
A_0)^{-1}(A_0 - A)}$ is differentiable almost everywhere.  Differentiating
the resolvent identity gives
$\frac{d}{d\theta}(e^{i\theta}I - A_0)^{-1} = -ie^{i\theta}
(e^{i\theta}I - A_0)^{-2}$, so
\[
  \left|\frac{d}{d\theta} m_{\mathrm{res}}(\theta)\right|
  \leq \norm{(e^{i\theta}I - A_0)^{-1}}^2 \cdot \norm{A_0 - A}.
\]
Setting $\kappa = \sup_{|z|=1}\norm{(zI-A_0)^{-1}}$, a valid Lipschitz bound is
$L = \kappa^2 \norm{A_0 - A}$.  Analytically, $\kappa$ can be bounded by the
pseudospectral radius or by the stability radius of $A_0$; numerically we
estimate $\kappa$ on the same contour grid.

%% ============================================================
\section{Reference Selection}
\label{sec:refsel}

The quality of both certificates depends critically on the choice of reference
$A_0$.  A reference that is far from $A$ in the resolvent norm will give a
large certificate margin only if $A_0$ is very stable (small resolvent norms),
and vice versa.

\subsection{Max-min formulation}

Given $A$ and a class $\mathcal{S}$ of Schur-stable matrices, define the
\emph{optimal reference} as
\[
  A_0^* = \argmax_{A_0 \in \mathcal{S}} \min_{\theta \in [0, 2\pi]}
  m(A, A_0, \theta),
\]
where $m$ is either the determinant or resolvent margin.  This is a max-min
(minimax) optimisation over a non-convex domain.

\subsection{Structured reference families}

We consider the following practical families $\mathcal{S}$:

\paragraph{Scalar.}
$A_0 = rI$ for scalar $r \in [0, 1)$.  Simple one-dimensional search over
$r$.  Worst-case for non-normal $A$ since the resolvent norm of $rI$ grows
as $1/(1-r)$.

\paragraph{Diagonal with shrunken spectrum.}
Given the eigenvalues $\lambda_i$ of $A$, set $a_{0,i} = s \min(|\lambda_i|, 1)
e^{i \arg(\lambda_i)}$ for a shrinkage factor $s \in (0, 1)$.  This
aligns phases with $A$ and shrinks radii to a stable baseline.

\paragraph{Random diagonal search.}
Sample random diagonal stable matrices $A_0 = \mathrm{diag}(r_i e^{i\phi_i})$
with $r_i \sim \mathrm{Uniform}(0, 0.95)$ and $\phi_i$ random, evaluate the
margin, keep the best.  Inexpensive and often effective for diagonal-structured $A$.

\paragraph{Future directions.}
Gradient-based optimisation of $\min_\theta m_{\mathrm{res}}(A, A_0, \theta)$
over the Lie group of stable matrices is a natural extension left to future
work (see also Paper~2 in this series).

%% ============================================================
\section{Experiments}
\label{sec:experiments}

\subsection{Setup}

All experiments use 64-bit complex arithmetic with the NumPy/SciPy stack.
Margins are evaluated on equispaced grids of $N = 512$ points unless stated
otherwise.  The resolvent certificate uses the matrix 2-norm.  We report both
the sampled-only verdict (positive minimum margin) and the rigorous verdict
(Theorem~\ref{thm:grid} with the numerical Lipschitz bound).

Unless otherwise noted, experiments use $n = 8$ (state dimension) and $M = 200$
random matrix instances per family.  Reference selection uses a random diagonal
search with 50 candidates followed by the shrunken-eigenvalue diagonal, with the
better result retained.

\subsection{Experiment 1: Certificate geometry}

We evaluate both certificates on the following matrix families:

\begin{description}
  \item[normal\_stable / normal\_unstable.] Eigenvalues uniformly distributed
    in a disk of radius 0.85 / 1.1, conjugated by a random unitary.
  \item[nonnormal\_stable / nonnormal\_unstable.] A stable / unstable
    non-normal matrix constructed by a random ill-conditioned similarity
    transform with departure 20; spectral radius rescaled to 0.85 / 1.05.
  \item[jordan\_stable / jordan\_near\_boundary.] Upper-triangular Jordan-like
    matrix with eigenvalue 0.80 / 0.95 and superdiagonal entries 0.5 / 1.0.
  \item[block\_diagonal\_stable.] Block-diagonal with $n/4 \times n/4$ random
    normal stable blocks.
  \item[dlr\_stable.] Diagonal-plus-low-rank with diagonal radius 0.85 and
    rank $\lfloor n/8 \rfloor$.
\end{description}

\TBD{Insert Table 1 here after running experiments/synthetic\_certification.py}

Key findings anticipated:
(i) resolvent certificates have uniformly lower false-negative rates than
    determinant certificates;
(ii) false-negative rates are highest for Jordan-near-boundary and non-normal
    matrices due to large resolvent norms;
(iii) reference selection (diagonal search) reduces false-negative rates by
    \TBD{X}\% over scalar references on non-normal families.

\subsection{Experiment 2: Discretisation correctness}

We verify Theorem~\ref{thm:grid} empirically.  For each test matrix we record
the sampled certificate, the rigorous certificate, and the ground-truth (dense
$N = 4096$ grid).  We check:
(a) no false positives for the rigorous certificate (by Theorem~\ref{thm:grid},
    none should occur);
(b) false positives for the sampled-only diagnostic at small $N$;
(c) convergence of the rigorous certificate as $N$ grows.

\TBD{Insert Table 2 / Figure 2 after running experiments/discretisation\_correctness.py}

\subsection{Experiment 3: Trained SSM case study}

We train small diagonal complex SSMs (state dimension $n = 8$) on
synthetic next-step prediction tasks from AR processes with poles at various
distances from the unit circle.  After training via L-BFGS-B, we evaluate
all three reference variants (scalar, shrunken-diagonal, random-search) and
report whether the resolvent certificate succeeds.

\TBD{Insert Table 3 / Figure 3 after running experiments/trained\_ssm.py}

\subsection{Experiment 4: Runtime and scalability}

We measure median wall-clock time for a single certificate evaluation as a
function of matrix dimension $n \in \{4, 8, 16, 32, 64, 128, 256\}$ and
grid size $K \in \{64, 128, 256, 512\}$, for dense normal matrices and for
diagonal-plus-low-rank matrices.

\TBD{Insert Table 4 / Figure 4 after running experiments/runtime\_scalability.py}

The resolvent certificate at grid size $K$ requires $K$ linear solves of
$n \times n$ systems, giving a dominant cost of $O(K n^3)$ for dense matrices.
For diagonal $A_0$ the linear solve reduces to $O(Kn)$, and for low-rank
corrections a Woodbury application adds $O(Kn r^2)$ where $r$ is the rank.

%% ============================================================
\section{Related Work}
\label{sec:related}

\paragraph{Stability of recurrent networks and SSMs.}
Spectral penalties \cite{sedghi2018characterizing}, Lyapunov regularisers,
and stable parameterisations~\cite{lezcano2019cheap} have all been proposed for
recurrent networks.  Our certificates are complementary: they are post-hoc
and do not modify training, but provide interpretable margin information.

\paragraph{Pseudospectra and stability radii.}
The stability radius of a matrix under unstructured perturbations equals the
smallest singular value of the resolvent on the unit circle, and is the
primary quantity bounded by our Theorem~\ref{thm:res}~\cite{hinrichsen1986stability,trefethen2005spectra}.
Our certificates provide a computable sufficient condition without requiring
the full pseudospectral computation.

\paragraph{Rouch\'e in control and signal processing.}
Rouch\'e's theorem has been applied to pole-placement certification in
robust control~\cite{bhattacharyya1995robust} and to zero-counting in filter
design.  Our contribution is its systematic application to learned transition
matrices in the SSM setting.

\paragraph{Finite-grid quadrature for contour integrals.}
The deterministic grid certificate (Theorem~\ref{thm:grid}) is related to
Euler--Maclaurin-style error bounds for periodic function quadrature.  The
Lipschitz-based bound we use is elementary; tighter trigonometric-polynomial
bounds are possible but unnecessary for the verification purpose here.

%% ============================================================
\section{Limitations}
\label{sec:limitations}

\begin{enumerate}
  \item \emph{Sufficiency only.}  A matrix can be Schur stable yet fail both
    certificates: this happens whenever no close stable reference exists with
    a favourable resolvent norm.
  \item \emph{Non-normal sensitivity.}  For highly non-normal $A_0$, the
    resolvent norm $\kappa = \sup_{|z|=1}\norm{(zI-A_0)^{-1}}$ can be orders
    of magnitude larger than $1/(1-\varrho(A_0))$, inflating the Lipschitz
    bound and requiring very fine grids for rigorous certification.
  \item \emph{Determinant ill-conditioning.}  For $n \geq 20$, $|p_0(z)|$
    can approach double-precision overflow or underflow, making the
    determinant certificate unreliable without log-determinant stabilisation.
  \item \emph{Computational cost.}  Dense resolvent checks cost $O(Kn^3)$.
    At $n = 256$ and $K = 512$ this is around $8 \times 10^9$ flops, which
    is feasible on a workstation but slow compared to a single eigenvalue
    decomposition ($O(n^3)$).
  \item \emph{Post-hoc only.}  The certificates assess a fixed trained matrix.
    Incorporating them into the training loss (a contour barrier) is deferred
    to Paper~2 in this series.
  \item \emph{Numerical Lipschitz bound.}  The bound $L = \kappa^2 \norm{A_0 - A}$
    is computed numerically; for a fully validated certificate an analytic
    bound on $\kappa$ is required (e.g.\ via interval arithmetic or pseudospectral
    radius bounds).
\end{enumerate}

%% ============================================================
\section{Conclusion}
\label{sec:conclusion}

We have developed and analysed two Rouch\'e-type contour certificates for
Schur stability of learned state-space transition matrices.  The resolvent
certificate is the practically preferred form: it is numerically stable for
moderate dimensions, connects cleanly to the stability radius and small-gain
framework, admits a simple interpretation for diagonal SSMs, and is
significantly less conservative than the determinant certificate.

The deterministic finite-grid theorem converts sampled margin checks into
mathematical certificates under a Lipschitz condition that can be estimated
from the same contour grid used for margin evaluation.  Reference selection---
choosing a stable $A_0$ with small resolvent norm that is close to $A$---is
the dominant practical lever for reducing false-negative rates.

We view these results as a first step toward a broader framework in which
complex-analytic tools complement algebraic stability conditions for learned
dynamical systems.  Training-time analogues (contour barriers, complex-analytic
Lyapunov functions), structured SSM-specific certificates, and
pathwise robustness analysis are directions pursued in companion papers.

%% ============================================================
\bibliographystyle{plain}
\bibliography{refs}

% Placeholder bibliography entries (replace with refs.bib)
\begin{thebibliography}{99}

\bibitem{ahlfors1979complex}
L.~V. Ahlfors.
\newblock {\em Complex Analysis}, 3rd ed.
\newblock McGraw-Hill, 1979.

\bibitem{bhattacharyya1995robust}
S.~P. Bhattacharyya, H.~Chapellat, and L.~H. Keel.
\newblock {\em Robust Control: The Parametric Approach}.
\newblock Prentice-Hall, 1995.

\bibitem{desoer1975feedback}
C.~A. Desoer and M.~Vidyasagar.
\newblock {\em Feedback Systems: Input--Output Properties}.
\newblock Academic Press, 1975.

\bibitem{gu2022s4}
A.~Gu, K.~Goel, and C.~R\'e.
\newblock Efficiently modeling long sequences with structured state spaces.
\newblock In {\em ICLR}, 2022.

\bibitem{gu2023mamba}
A.~Gu and T.~Dao.
\newblock Mamba: Linear-time sequence modeling with selective state spaces.
\newblock {\em arXiv:2312.00752}, 2023.

\bibitem{hinrichsen1986stability}
D.~Hinrichsen and A.~J. Pritchard.
\newblock Stability radii of linear systems.
\newblock {\em Systems \& Control Letters}, 7(1):1--10, 1986.

\bibitem{jury1964inners}
E.~I. Jury.
\newblock {\em Theory and Application of the z-Transform Method}.
\newblock Wiley, 1964.

\bibitem{lezcano2019cheap}
M.~Lezcano-Casado and D.~Mart\'inez-Rubio.
\newblock Cheap orthogonal constraints in neural networks.
\newblock In {\em ICML}, 2019.

\bibitem{sedghi2018characterizing}
H.~Sedghi and A.~Anandkumar.
\newblock Provable tensor methods for learning mixtures of classifiers.
\newblock {\em arXiv:1412.3046}, 2014.

\bibitem{smith2023lru}
J.~T.~H. Smith, A.~Warrington, and S.~W. Linderman.
\newblock Simplified state space layers for sequence modeling.
\newblock In {\em ICLR}, 2023.

\bibitem{trefethen2005spectra}
L.~N. Trefethen and M.~Embree.
\newblock {\em Spectra and Pseudospectra}.
\newblock Princeton University Press, 2005.

\end{thebibliography}

\end{document}
